%
% reversi_120_appendix_regole.tex
%
% This file is part of the reversi program
% http://github.com/rcrr/reversi-v2
%
% Author Roberto Corradini mailto:rob_corradini@yahoo.it
% Copyright 2026 Roberto Corradini. All rights reserved.
%
% License
%
% This program is free software; you can redistribute it and/or modify it
% under the terms of the GNU General Public License as published by the
% Free Software Foundation; either version 3, or (at your option) any
% later version.
%
% This program is distributed in the hope that it will be useful,
% but WITHOUT ANY WARRANTY; without even the implied warranty of
% MERCHANTABILITY or FITNESS FOR A PARTICULAR PURPOSE. See the
% GNU General Public License for more details.
%
% You should have received a copy of the GNU General Public License
% along with this program; if not, write to the Free Software
% Foundation, Inc., 59 Temple Place - Suite 330, Boston, MA  02111-1307, USA
% or visit the site <http://www.gnu.org/licenses/>.
%

%-----------------------------------------------------------------------------
% Appendice B — Regole del gioco.
%
% Le regole che seguono sono tradotte dalla pubblicazione della World Othello
% Federation (WOF), la fonte autorizzata delle regole ufficiali del gioco. Le
% figure sono ricostruite con il package othelloboard a partire dai diagrammi
% originali della pagina WOF. La fonte e' citata espressamente.
%-----------------------------------------------------------------------------

\captionsetup[subfigure]{justification=centering}

%--- Helpers per le figure delle regole ---------------------------------------
% \rulesboard{<stringa 64 char, UNA sola riga>}{<annotazioni e frecce>}
%   Disegna una scacchiera 8x8 (con i riferimenti a-h / 1-8) a partire da una
%   stringa di 64 caratteri (O = disco bianco, X = disco nero, - = cella vuota),
%   letti riga per riga (riga 1 in alto, da colonna a a h), e applica le
%   annotazioni/frecce passate come secondo argomento.
\newcommand{\rulesboard}[2]{%
  \begin{othelloboard}{1}%
    \drawboardfromstring{#1}%
    #2%
  \end{othelloboard}%
}

% Frecce ROSSE: visibili sia sui dischi neri sia sulle caselle vuote.
\newcommand{\arrR}[2]{\put(#1,#2){\textcolor{red}{\vector(1,0){14}}}}
\newcommand{\arrL}[2]{\put(#1,#2){\textcolor{red}{\vector(-1,0){14}}}}
\newcommand{\arrU}[2]{\put(#1,#2){\textcolor{red}{\vector(0,1){14}}}}
\newcommand{\arrD}[2]{\put(#1,#2){\textcolor{red}{\vector(0,-1){14}}}}

\chapter{Regole del gioco}

Questa appendice riporta le regole ufficiali del gioco, tradotte dalla
pubblicazione della \textbf{World Othello Federation (WOF)}:

\begin{quote}\small
\emph{Official rules for the game Othello}, World Othello Federation.\\
\url{https://www.worldothello.org/about/about-othello/othello-rules/official-rules/english}
\end{quote}

\section{Obiettivo del gioco}

L'obiettivo del gioco è avere la maggioranza dei dischi rivolti verso l'alto
sulla scacchiera, del proprio colore, alla fine della partita.

\section{Un minuto per impararlo\ldots}

Ogni giocatore prende 32 dischi e sceglie un colore da usare per tutta la
partita. Nero piazza due dischi neri e Bianco piazza due dischi bianchi come
mostrato in Figura~\ref{fig:rules-setup}.

\begin{figure}[htbp]
  \centering
  \scalebox{0.62}{%
    \rulesboard{---------------------------OX------XO---------------------------}{}%
  }
  \caption{La partita inizia sempre con questa disposizione.}
  \label{fig:rules-setup}
\end{figure}

Una mossa consiste nel ``aggirare'' il disco (o i dischi)
dell'avversario, per poi ribaltare il disco (o i dischi) aggirato al
proprio colore. Aggirare significa piazzare un disco sulla scacchiera
in modo che la fila (o le file) di dischi dell'avversario sia delimitata a
ciascuna estremità da un disco del proprio colore. (Una ``fila'' può essere
composta da uno o più dischi).

\begin{quote}\small
\textbf{Il termine.} Nel linguaggio regolamentare l'operazione si chiama
\emph{aggiramento} (verbo: \emph{aggirare}): ``una mossa è valida se permette
l'aggiramento di una o più pedine avversarie.'' Nel gergo tra giocatori si
usano anche espressioni come \emph{intrappolare}, \emph{chiudere a panino},
\emph{prendere in mezzo} o \emph{stringere}, a indicare la stessa idea:
mettere i dischi avversari in mezzo a due dischi del proprio colore.
\end{quote}

Ecco un esempio. Il disco bianco A era già a posto sulla scacchiera. La
collocazione del disco bianco B aggira la fila di tre dischi neri.
Bianco ribalta i dischi aggirati e ora la fila è così:

\begin{figure}[htbp]
  \centering
  \begin{subfigure}[b]{0.45\textwidth}
    \centering
    \scalebox{0.5}{%
      \rulesboard{--------------------------OXXXO---------------------------------}%
        {\posannotation{C4}{A}\posannotation{G4}{B}}%
    }
    \caption{La mossa di aggiramento: viene posto il disco bianco B.}
  \end{subfigure}
  \hfill
  \begin{subfigure}[b]{0.45\textwidth}
    \centering
    \scalebox{0.5}{%
      \rulesboard{--------------------------OOOOO---------------------------------}%
        {\posannotation{C4}{A}\posannotation{G4}{B}}%
    }
    \caption{I tre dischi aggirati vengono ribaltati in bianco.}
  \end{subfigure}
  \caption{Esempio di aggiramento di una fila di tre dischi.}
  \label{fig:rules-outflank}
\end{figure}

\section{Regole di Othello}

\begin{enumerate}
  \item Nero muove sempre per primo.

  \item Se al tuo turno non puoi aggirare e ribaltare almeno un disco
    avversario, il tuo turno è annullato e l'avversario muove di nuovo.
    Tuttavia, se una mossa è a tua disposizione, non puoi cedere il turno.
    % L'elenco originale non ha una regola 3; la numerazione riprende da 4.
    \setcounter{enumi}{3}

    \begin{figure}[htbp]
      \centering
      \begin{subfigure}[b]{0.45\textwidth}
        \centering
        \scalebox{0.5}{%
          \rulesboard{----------------XO--------O-------O-------X---------------------}%
            {\arrU{60}{118}}%
        }
        \caption{Disco posto}
      \end{subfigure}
      \hfill
      \begin{subfigure}[b]{0.45\textwidth}
        \centering
        \scalebox{0.5}{%
          \rulesboard{----------------XXX-------X-------X-------X---------------------}%
            {\arrR{22}{132}\arrU{60}{94}\arrU{60}{70}}%
        }
        \caption{Questi dischi ribaltati}
      \end{subfigure}
      \caption{La mossa (a) e i dischi ribaltati (b).}
      \label{fig:rules-r2}
    \end{figure}

  \item I giocatori non possono saltare oltre i propri dischi per aggirare
    un disco avversario. (Vedi Figura~\ref{fig:rules-r4}.)

    \begin{figure}[htbp]
      \centering
      \scalebox{0.62}{%
        \rulesboard{------------------------------------------------XOXOOX----------}%
          {\arrD{12}{50}\posannotation{B7}{1}}%
      }
      \caption{Questo disco aggira e ribalta solo il disco bianco 1.}
      \label{fig:rules-r4}
    \end{figure}

  \item I dischi possono essere aggirati solo come risultato diretto di
    una mossa e devono trovarsi nella linea diretta del disco posto.
    (Vedi Figure~\ref{fig:rules-r5}a e~\ref{fig:rules-r5}b.)

    \begin{figure}[htbp]
      \centering
      \begin{subfigure}[b]{0.45\textwidth}
        \centering
        \scalebox{0.5}{%
          \rulesboard{X-------O-------OOOX----X---------------------------------------}%
            {\arrR{-2}{180}}%
        }
        \caption{Disco posto qui}
      \end{subfigure}
      \hfill
      \begin{subfigure}[b]{0.45\textwidth}
        \centering
        \scalebox{0.5}{%
          \rulesboard{X-------X-------XOOX----X---------------------------------------}%
            {\arrR{-2}{156}\arrR{-2}{132}\posannotation{B3}{1}\posannotation{C3}{2}}%
        }
        \caption{Questi 2 dischi ribaltati. I dischi 1 e 2 non vengono ribaltati (anche se sembrano aggirati).}
      \end{subfigure}
      \caption{Vengono ribaltati solo i dischi nella linea diretta della mossa.}
      \label{fig:rules-r5}
    \end{figure}

  \item Tutti i dischi aggirati in una mossa devono essere ribaltati,
    anche se sarebbe a vantaggio del giocatore non ribaltarli affatto.

  \item Una volta che un disco è posto su una casella, non può più essere
    mosso su un'altra casella più tardi nella partita.

  \item Quando non è più possibile per nessuno dei due giocatori muovere, la
    partita è finita. I dischi vengono contati e il giocatore con la
    maggioranza del proprio colore visibile è il vincitore.

    \begin{quote}\itshape
      Nota: è possibile che una partita finisca prima che tutte e 64 le
      caselle siano riempite.
    \end{quote}
\end{enumerate}

\section{Partita di esempio}

Nero muove per primo. Un disco nero può essere posto in C4, D3, E6 o F5 per
aggirare un disco bianco. Nero decide di posizionare un disco in C4
(Vedi Figura~\ref{fig:rules-sg}a). Il disco bianco aggirato tra i due
dischi neri viene ribaltato (Vedi Figura~\ref{fig:rules-sg}b). Ora è il turno
di Bianco. Un disco bianco può essere posto in C3, E3 o C5 per aggirare
un disco nero. Bianco decide di posizionare un disco in C3 (Vedi
Figura~\ref{fig:rules-sg}c). Il disco nero aggirato tra i due dischi
bianchi viene ribaltato (Vedi Figura~\ref{fig:rules-sg}d).

\begin{figure}[htbp]
  \centering
  \begin{subfigure}[b]{0.45\textwidth}
    \centering
    \scalebox{0.5}{%
      \rulesboard{--------------------------XOX------XO---------------------------}{}%
    }
    \caption{Nero posiziona un disco in C4.}
  \end{subfigure}
  \hfill
  \begin{subfigure}[b]{0.45\textwidth}
    \centering
    \scalebox{0.5}{%
      \rulesboard{--------------------------XXX------XO---------------------------}{}%
    }
    \caption{Il disco bianco aggirato (D4) viene ribaltato.}
  \end{subfigure}

  \begin{subfigure}[b]{0.45\textwidth}
    \centering
    \scalebox{0.5}{%
      \rulesboard{------------------O-------XXX------XO---------------------------}{}%
    }
    \caption{Bianco posiziona un disco in C3.}
  \end{subfigure}
  \hfill
  \begin{subfigure}[b]{0.45\textwidth}
    \centering
    \scalebox{0.5}{%
      \rulesboard{------------------O-------XOX------XO---------------------------}{}%
    }
    \caption{Il disco nero aggirato (D4) viene ribaltato.}
  \end{subfigure}
  \caption{Le prime due coppie di mosse della partita di esempio.}
  \label{fig:rules-sg}
\end{figure}

La partita prosegue in questo modo fino a quando nessuno dei due giocatori può
più muovere. Il giocatore con più dischi a questo punto è il vincitore.

\section{Punteggio e fine anticipata della partita}

Nei campionati internazionali regolati dalla World Othello Federation (WOF),
quando la partita termina in anticipo — perché nessun giocatore ha più mosse
valide, oppure perché uno dei due subisce un ``cappotto'' (\emph{wipeout})
perdendo tutte le proprie pedine — le caselle rimaste vuote vengono assegnate
d'ufficio al vincitore. La somma dei punti dei due giocatori nei verbali
ufficiali deve quindi essere sempre pari a 64: la differenza finale include
anche tutte le caselle vuote accumulate dal vincitore.

\begin{quote}\small
\emph{World Othello Championships Rules}, Sezione IV.8 ``Scoring'':\\
``The official score of the game will be determined by counting up the discs
of each colour on the board, counting empty squares for the winner.''\\
(Traduzione: ``Il punteggio ufficiale della partita sarà determinato contando
i dischi di ciascun colore sulla scacchiera, contando le caselle vuote a
favore del vincitore.'')\\
\url{https://www.worldothello.org/application/files/8816/9235/2069/NEW_WOC_Rules_2023.pdf}
\end{quote}
